\documentclass[11pt]{article}
\usepackage[margin=1in]{geometry}
\usepackage[T1]{fontenc}
\usepackage{lmodern,amsmath,amsfonts,booktabs,graphicx}
\usepackage[hidelinks]{hyperref}
\usepackage{xurl}
\setlength{\parindent}{0pt}
\setlength{\parskip}{6pt}
\setlength{\emergencystretch}{3em}
\graphicspath{{simulation_output/estimator_near_visibility_20261003/}}
\title{Why the Nearer Recording Retains More Position Bias}
\author{Pan Dynamics Collision Avoidance Research}
\date{October 3, 2026}
\begin{document}
\maketitle

\section{Finding: visibility loss and geometric bias}
The closer recording has more points but loses low vehicle surfaces through
the LiDAR's downward field-of-view limit and the camera's image boundary.
The current fitter's ``lower body'' is defined relative to the height interval
of the \emph{visible} cluster. Once the actual lower body is missing, this
rule retains a higher, differently shaped band. Its fitted rectangle can
therefore have a consistent longitudinal offset even with dense returns.

This explanation is supported by a controlled point-cloud experiment:
five existing farther clusters were translated to the same near distance,
with their shape, noise and initial fit pose held fixed across four masks.
Applying only the sensor visibility limits increases longitudinal bias in
all five. This identifies a concrete visibility--fitting mechanism without
claiming an exact causal percentage of the recorded error.

All performance statistics here use the previously declared $t\geq2$ s
window, frames 100--299 of the normal recordings. Startup error is outside
the question being addressed. Production source, estimator gains and
initialization were not changed. No live CARLA connection, new capture,
detector inference or MATLAB replay was required.

\section{Actual distances and bias in the evaluation window}
The labels 15 m and 20 m describe initial scenario spacing. They do not
describe the distance during evaluation:
\begin{center}\small
\begin{tabular}{@{}lrr@{}}\toprule
Quantity after 2 s & 15 m recording & 20 m recording \\\midrule
Target center longitudinal distance (m) & 4.729--4.812 & 4.930--9.048 \\
Raw position RMSE (m) & 0.12458 & 0.04324 \\
Mean raw longitudinal error (m) & +0.12378 & +0.03809 \\
Raw longitudinal standard deviation (m) & 0.00455 & 0.00710 \\
Mean raw lateral error (m) & -0.00127 & -0.00402 \\
Centered 2D raw-error RMS (m) & 0.01403 & 0.02007 \\
Squared mean bias / position MSE & 98.73\% & 78.46\% \\
Image boxes touching bottom ($y_{\max}\geq639$) & 200/200 & 32/200 \\
Estimated position RMSE (m) & 0.15212 & 0.05221 \\
\bottomrule\end{tabular}
\end{center}
The nearer recording actually has smaller longitudinal fluctuations, but
its approximately 12.4 cm mean offset dominates RMSE. The decomposition
$\mathbb E\|e\|^2=\|\mathbb E e\|^2+
\mathbb E\|e-\mathbb E e\|^2$ is an error accounting identity, not a causal
allocation. Centered variation includes scene variation and fitting effects;
it is not a direct measurement of LiDAR noise alone. Both recordings specify
\texttt{noise\_stddev}=0.02 m, rather than different noise settings.

\section{Why near range removes low surfaces}
The recorded LiDAR origin is approximately $(x,z)=(-0.15,1.8)$ m in ego
coordinates, with lowest elevation $-30^\circ$. The front camera origin is
$(1.55,1.4)$ m with a 640 by 640 image and $100^\circ$ horizontal field of
view. Its square image gives the same vertical angular extent. These are
ego-coordinate heights, not heights above the road.

At frame 109, the target center is $x=4.79886$ m in the 15 m recording.
The declared 5 m rectangle puts its rear midpoint approximately 2.44890 m
horizontally from the LiDAR. At this rear location the lower LiDAR ray is
\[
 z_{\min,L}=1.8-2.44890\tan30^\circ=0.38612\ \mathrm{m}.
\]
At the same declared rear plane the lower image boundary is about
$z_{\min,C}=0.50754$ m. Thus low returns can be missing either because the
LiDAR never samples that direction or because their image projections fall
below the camera. These values describe a particular declared rear plane;
the cutoffs vary with each point's position, so they are not global constant
height thresholds for the cloud.

The measured height distribution confirms a different visible band:
\begin{center}\small
\begin{tabular}{@{}lrr@{}}\toprule
Frame 109 & 15 m recording & 20 m recording \\\midrule
Target center $x$ (m) & 4.79886 & 8.63915 \\
Associated points & 999 & 269 \\
5th-percentile point height (m) & 0.39895 & 0.04602 \\
Current lower-body selection ceiling (m) & 0.74543 & 0.58689 \\
Raw longitudinal bias (m) & +0.11245 & +0.03768 \\
\bottomrule\end{tabular}
\end{center}
In the nearer cloud, the lower half of the visible height interval is already
well above the low body visible in the farther cloud. Increasing point count
within this truncated band does not recreate the missing footprint support.
Near the end of the 20 m recording, its center distance also decreases to
4.93029 m: the height 5th percentile rises to 0.32040 m and raw longitudinal
bias rises to 0.07256 m. This within-record trend is consistent with the
visibility explanation, but by itself is observational evidence.

\section{Same-cloud visibility controls}
Five native clusters from 20 m recording frames 100, 109, 121, 150 and 200
were translated along ego $x$ to a common target-center $x=4.798859$ m.
The following were held constant within each four-way comparison: all original
XYZ returns before masking, their noise, declared dimensions, production
fitting code and a common truth-pose numerical seed. The four masks were
no crop, LiDAR elevation limits only, full camera-image bounds only, and both.
Truth is used only for this explicitly diagnostic construction and scoring.
It is not supplied to deployed perception.

\begin{center}\small
\begin{tabular}{@{}rrrrr@{}}\toprule
Source frame & No crop & LiDAR limits & Camera bounds & Both \\
 & \multicolumn{4}{c}{Fitted longitudinal bias (cm)} \\\midrule
100 & 3.81 & 9.45 & 10.28 & 10.28 \\
109 & 3.70 & 9.00 & 10.01 & 10.00 \\
121 & 3.62 & 7.74 & 9.93 & 9.93 \\
150 & 3.18 & 10.10 & 11.19 & 11.19 \\
200 & 3.62 & 9.74 & 10.72 & 10.72 \\
\bottomrule\end{tabular}
\end{center}
For source frame 109, the original 269 points become 202 with the LiDAR mask,
189 with the camera mask and 188 with both. The masks overlap strongly, so
their effects cannot be added. The resulting approximately 10 cm error is
close in scale to the actual near frame's 11.2 cm longitudinal error, but
this is not a complete simulation of that frame.

The experiment applies geometric visibility to already recorded points.
It does not rerun ray casting, occlusion, discrete LiDAR beam sampling,
noise generation, YOLO selection, vehicle pitch/roll or estimator dynamics
at the translated pose. Twenty controlled fits completed; five saved-input
fit reproductions agree within $8.9\times10^{-12}$ m. The comparison supports
the mechanism while leaving shape and dimension contributions unseparated.

\begin{figure}[p]
\centering\includegraphics[width=\textwidth]{visibility_controls.pdf}
\caption{Left: a translated actual cluster loses its low returns under the
recorded sensor visibility constraints. The blue ray is the $y=0$ LiDAR
section; the actual mask uses every point's horizontal radius. Right: the
same five clouds acquire larger longitudinal bias when only visibility
cropping changes. This is an offline diagnostic, not newly captured data.}
\end{figure}

\section{Implication for the current algorithm}
The immediate limitation is a visible-surface-to-footprint model mismatch.
The fitter places an approximate rectangle against the current visible body
band; an inward-recessed band shifts the inferred rear plane forward and,
through the half-length offset, shifts the center forward. The observed
positive longitudinal bias and controlled cropping experiment support this
mechanism. The current estimator receives that biased position stream and
does not have an independent measurement that identifies the missing offset.

Further steady-tracking improvement should explicitly handle truncated body
visibility and use measured support that represents the footprint boundary,
or model the relation between visible surfaces and that boundary. Changing
observer initialization is outside this diagnosis. A fixed 12 cm subtraction
would not address the demonstrated dependence on visibility and range.

The analyzed perception commit is
\path{21c7a0c2d4bc95dfe9d43a8ada5f48d69d64a8b4}.
Compact measurements, all 20 controls, source/input hashes and validation
details accompany this report in
\path{report/ESTIMATOR_NEAR_VISIBILITY_20261003/}.
The executable local diagnosis and generated figure remain in
\path{simulation_output/estimator_near_visibility_20261003/}.
Production algorithms and all concurrent controller, estimator, adapter,
theory and test edits are outside this report-only change. Existing frozen
recordings were inspected; no new unit suite was necessary or executed.
\end{document}
